\section{SmallWood Proof-Stack Details}
\label{app:smallwood}

This appendix records schematic compiled surfaces used by the issuance and showing statements.  It is descriptive:
the semantic statements are those of Section~\ref{sec:smallwood-model}, and final proof-size claims require the
complete row inventory, constraint degrees, and soundness calculation for the instantiated relation.

\subsection{Row inventories}
\paragraph{Issuance rows.}
The issuance proof commits rows encoding:
\[
  M,
  \qquad
  \vecm,
  \qquad
  \veck,
  \qquad
  s,
  \qquad
  e.
\]
The active relations are
\begin{align*}
  c&=C_MM+A_s s+e,\\
  M&=\vecm\Vert\veck,\\
  \veck&\in\mathcal K_{\mathsf{key}},
\end{align*}
plus the prescribed coefficient bounds and issuance policy constraints.  In the bounded-key instantiation, key
membership is implemented with the same signed-lift bounded-membership convention as the attribute rows, together
with any additional key-space equations.

\paragraph{Showing rows.}
The showing proof commits rows encoding:
\[
  \vecu,
  \quad
  M,
  \quad
  \vecm,
  \quad
  \veck,
  \quad
  s,
  \quad
  e,
  \quad
  \muSig,
  \quad
  x_0,
  \quad
  x_1,
  \quad
  Z,
\]
plus rows for the signature shortness gadget and the PRF companion relation.  The active algebraic relations are
\begin{align*}
  (B_3-\mathbf 1_nx_1)\Had Z&=1_n,\\
  \mA\vecu&=B_0+B_1\muSig+B_2x_0+Z+C_MM+A_s s+e,\\
  M&=\vecm\Vert\veck,\\
  \prftag&=\PRF(\veck,\nonce).
\end{align*}
Equivalently, the showing proof implies the eliminated coefficient-domain relation
\[
  (B_3-\mathbf 1_nx_1)\Had
  \bigl(\mA\vecu-B_0-B_1\muSig-B_2x_0-C_MM-A_s s-e\bigr)=1_n.
\]

\subsection{Commitment constraints}
The commitment opening constraint is linear over \(\Rq\) and is compiled coefficient-wise into SmallWood rows:
\[
  c=C_MM+A_s s+e.
\]
The block \(I_{n_c}\) multiplying \(e\) is public and requires no invertibility check.  Binding of two distinct
openings is external to the proof system and follows from the Module-SIS assumption for
\([C_M\mid A_s\mid I_{n_c}]\).

\subsection{BB-tran constraints}
The BB-tran constraints are expressed in witness form.  Showing uses the inverse equation
\[
  (B_3-\mathbf 1_nx_1)\Had Z=1_n
\]
and the signature surface
\[
  \mA\vecu=B_0+B_1\muSig+B_2x_0+Z+C_MM+A_s s+e.
\]
The first equation is bilinear in \((x_1,Z)\).  The second equation is linear in the committed witness rows once
\(Z\) is carried explicitly.

\subsection{PRF companion constraints}
The PRF companion relation binds the public pair \((\nonce,\prftag)\) to the hidden key \(\veck\) carried inside
\(M\).  The proof commits key-lane rows, selected S-box checkpoint rows, and any auxiliary rows needed to replay the
Poseidon2-style linear layers between checkpoints.  Selector constraints tie the key-lane rows to the signed key
coordinates and enforce \(\veck\in\mathcal K_{\mathsf{key}}\); final feed-forward and truncation constraints tie the
public tag to the resulting PRF output.  Security of this keyed construction is the explicit PRF assumption of
Assumption~\ref{ass:poseidon-prf}.

\subsection{Single committed witness}
All rows in one proof live under one SmallWood commitment.  This ensures that the prover cannot use one hidden
message for the commitment term, a different rational-hash characteristic for the signature relation, and a third key
for the PRF.  In the showing proof, the same hidden key \(\veck\) is used simultaneously in the semantic message
encoding and the tag relation.
